\documentclass[10pt,aspectratio=169]{beamer}
\usepackage{bma02}

\title[Longitud de arco]{Longitud de arco}
\subtitle{Curvas paramétricas y curvas polares\\[2pt]\small Semana 5 -- Sesión 2}

\begin{document}

\begin{frame}
  \titlepage
\end{frame}

\begin{frame}{Objetivos de la sesión}
  \begin{enumerate}
    \item Definir la \resaltar{longitud de una curva} como límite de longitudes de poligonales inscritas.
    \item Deducir y aplicar la fórmula $L=\int\sqrt{x'^2+y'^2}\dt$ para curvas paramétricas.
    \item Obtener los casos particulares: gráficas $y=f(x)$ y \resaltar{curvas polares} $r=f(\theta)$.
    \item Introducir la función longitud de arco $s(t)$.
  \end{enumerate}
  \vspace{1mm}
  {\small\textbf{Competencia 2 del sílabo:} Describe procesos naturales y físicos con modelos matemáticos que usan la integración.}
  \begin{block}{Plan (2 horas; incluye verificación previa de 5 min)}
    \begin{tabular}{ll}
      0:00 -- 0:25 & Definición: poligonales inscritas. Curvas rectificables\\
      0:25 -- 1:00 & Fórmula para curvas paramétricas y gráficas. Ejemplos\\
      1:00 -- 1:40 & Longitud de arco de curvas polares\\
      1:40 -- 2:00 & Actividad colaborativa, resumen y cierre metacognitivo
    \end{tabular}
  \end{block}
\end{frame}

\begin{frame}{Aula invertida: preparación y verificación}
  \small
  \begin{previa}[antes de la sesión]
    \begin{enumerate}
      \item Lea la definición de longitud de arco con poligonales inscritas y las fórmulas.
      \item Escriba la longitud de $y=\frac23x^{3/2}$ en $[0,3]$ como integral.
      \item Pregunta guía: ¿por qué la longitud es el límite de longitudes de poligonales?
    \end{enumerate}
  \end{previa}
  \begin{block}{Verificación al inicio (5 min, en pares)}
    \begin{enumerate}
      \item Escriba la fórmula de longitud de $y=f(x)$.
      \item Escriba la fórmula de longitud en polares.
    \end{enumerate}
    Los pares discuten y votan; el docente aclara solo lo que genera dudas y ajusta el ritmo de la sesión.
  \end{block}
\end{frame}

\section{Definición de longitud de arco}

\begin{frame}{Aproximación por poligonales}
  \begin{columns}
    \begin{column}{0.5\textwidth}
      Sea $C:\ \mathbf r(t)=(x(t),y(t))$, $t\in[\alpha,\beta]$, y $P=\{t_0,\dots,t_n\}$ una partición. La poligonal
      inscrita une $P_i=\mathbf r(t_i)$:
      \[ \ell(P)=\sum_{i=1}^n|P_{i-1}P_i|. \]
      \begin{definicion}
        $C$ es \resaltar{rectificable} si $\{\ell(P)\}$ está acotado; su \resaltar{longitud} es
        \[ L(C)=\sup_P\ell(P). \]
      \end{definicion}
      \small (Refinar una partición no disminuye $\ell(P)$, por la desigualdad triangular.)
    \end{column}
    \begin{column}{0.5\textwidth}
      \centering
      \begin{tikzpicture}
        \begin{axis}[bma, xmin=0, xmax=4.2, ymin=0, ymax=3, width=7cm, height=5.6cm, xtick=\empty, ytick=\empty, domain=0.3:3.9]
          \addplot[curva, samples=100] {1.3+0.9*sin(deg(1.6*x))+0.15*x};
          \addplot[uniGranate, thick, mark=*, mark size=1.5pt, samples=5] {1.3+0.9*sin(deg(1.6*x))+0.15*x};
          \addplot[uniVerde, thick, dashed, mark=*, mark size=1pt, samples=9] {1.3+0.9*sin(deg(1.6*x))+0.15*x};
          \node[uniGranate] at (axis cs:3.2,0.4) {\scriptsize $n=4$};
          \node[uniVerde] at (axis cs:3.2,0.15) {\scriptsize $n=8$};
          \node[uniGranate] at (axis cs:0.3,0.95) {\scriptsize $P_0$};
          \node[uniGranate] at (axis cs:3.9,2.2) {\scriptsize $P_n$};
        \end{axis}
      \end{tikzpicture}
    \end{column}
  \end{columns}
\end{frame}

\begin{frame}{Fórmula de la longitud de arco}
  \begin{teorema}
    Si $x',y'$ son continuas en $[\alpha,\beta]$ y $C$ se recorre una sola vez, entonces $C$ es rectificable y
    \[ L=\int_\alpha^\beta\sqrt{x'(t)^2+y'(t)^2}\dt. \]
  \end{teorema}
  \textbf{Idea de la prueba.} Por el TVM, $x(t_i)-x(t_{i-1})=x'(\xi_i)\Delta t_i$ y $y(t_i)-y(t_{i-1})=y'(\eta_i)\Delta t_i$:
  \[ |P_{i-1}P_i|=\sqrt{x'(\xi_i)^2+y'(\eta_i)^2}\,\Delta t_i. \]
  Es ``casi'' una suma de Riemann de $\sqrt{x'^2+y'^2}$ (los puntos $\xi_i,\eta_i$ pueden diferir; la continuidad
  uniforme controla esa diferencia).
  \begin{block}{Notación diferencial}
    \[ ds=\sqrt{dx^2+dy^2},\qquad L=\int_C ds. \]
  \end{block}
\end{frame}

\begin{frame}{Caso particular: gráficas de funciones}
  \begin{columns}
    \begin{column}{0.55\textwidth}
      Si $y=f(x)$, $a\le x\le b$, con $f'$ continua, tomamos $x=t$, $y=f(t)$:
      \begin{teorema}
        \[ L=\int_a^b\sqrt{1+f'(x)^2}\dx. \]
        Análogamente, si $x=g(y)$, $c\le y\le d$: \ $L=\displaystyle\int_c^d\sqrt{1+g'(y)^2}\dy$.
      \end{teorema}
      \begin{ejemplo}[1]
        $y=x^{3/2}$, $0\le x\le4$: \ $f'(x)=\frac32x^{1/2}$,
        \[ L=\int_0^4\sqrt{1+\tfrac94x}\dx=\frac{8}{27}\Big[(1+\tfrac94x)^{3/2}\Big]_0^4=\frac{8}{27}\big(10\sqrt{10}-1\big). \]
      \end{ejemplo}
    \end{column}
    \begin{column}{0.45\textwidth}
      \centering
      \begin{tikzpicture}
        \begin{axis}[bma, xmin=0, xmax=4.5, ymin=0, ymax=8.6, width=6cm, height=5.6cm, xtick={4}, ytick={8}, domain=0:4]
          \addplot[curva] {x^1.5} node[pos=0.9,left] {$y=x^{3/2}$};
          \addplot[uniGranate, thick, dashed] coordinates {(0,0) (4,8)};
        \end{axis}
      \end{tikzpicture}

      \small $L\approx9.07$; \ la cuerda mide $\sqrt{80}\approx8.94$.
    \end{column}
  \end{columns}
\end{frame}

\begin{frame}{Más ejemplos con gráficas}
  \begin{ejemplo}[2]
    $y=\ln(\cos x)$, $0\le x\le\frac\pi4$: \ $1+f'^2=1+\tan^2x=\sec^2x$,
    \[ L=\int_0^{\pi/4}\sec x\dx=\Big[\ln|\sec x+\tan x|\Big]_0^{\pi/4}=\ln\big(\sqrt2+1\big). \]
  \end{ejemplo}
  \begin{ejemplo}[3: catenaria]
    $y=\cosh x$, $0\le x\le a$: \ $1+\senh^2x=\cosh^2x$,
    \quad $L=\displaystyle\int_0^a\cosh x\dx=\senh a$. \ (Forma de un cable colgante.)
  \end{ejemplo}
  \begin{observacion}
    La raíz $\sqrt{1+f'^2}$ rara vez tiene primitiva elemental. Los ejemplos ``calculables'' se diseñan para que
    $1+f'^2$ sea un cuadrado perfecto. En general se usa integración numérica (Semana 9).
  \end{observacion}
\end{frame}

\section{Curvas paramétricas}

\begin{frame}{Ejemplos con curvas paramétricas}
  \begin{columns}
    \begin{column}{0.6\textwidth}
      \begin{ejemplo}[4: circunferencia]
        $x=r\cos t$, $y=r\sen t$: \ $L=\displaystyle\int_0^{2\pi}r\dt=2\pi r$. \quad ¡La definición de $\pi$ es coherente!
      \end{ejemplo}
      \begin{ejemplo}[5: astroide]
        $x=a\cos^3t$, $y=a\sen^3t$: \ $x'^2+y'^2=9a^2\sen^2t\cos^2t$.
        Por simetría (en el primer cuadrante $|\sen t\cos t|=\sen t\cos t$):
        \[ L=4\int_0^{\pi/2}3a\sen t\cos t\dt=6a. \]
      \end{ejemplo}
      \begin{observacion}
        En $[0,2\pi]$, $\int_0^{2\pi}3a\sen t\cos t\dt=0$: ¡no olvide el valor absoluto!
      \end{observacion}
    \end{column}
    \begin{column}{0.4\textwidth}
      \centering
      \begin{tikzpicture}
        \begin{axis}[bma, xmin=-1.3, xmax=1.3, ymin=-1.3, ymax=1.3, axis equal image, width=5.5cm, height=5.5cm, xtick={1}, xticklabels={$a$}, ytick={1}, yticklabels={$a$}]
          \addplot[uniVerde, very thick, domain=0:360, samples=200] ({cos(x)^3},{sin(x)^3});
          \addplot[uniGranate, ultra thick, domain=0:90, samples=60] ({cos(x)^3},{sin(x)^3});
        \end{axis}
      \end{tikzpicture}

      \small Un cuarto mide $\frac32a$.
    \end{column}
  \end{columns}
\end{frame}

\begin{frame}{Ejemplo 6: la elipse y las integrales elípticas}
  \begin{columns}
    \begin{column}{0.58\textwidth}
      $x=a\cos t$, $y=b\sen t$ con $a>b$. Sea $e=\sqrt{1-b^2/a^2}$ la excentricidad.
      \[ L=4\int_0^{\pi/2}\sqrt{a^2\sen^2t+b^2\cos^2t}\dt=4a\int_0^{\pi/2}\sqrt{1-e^2\cos^2t}\dt. \]
      Esta es una \resaltar{integral elíptica}: \emph{no} tiene primitiva elemental.
      \begin{observacion}[aproximación de Ramanujan]
        $L\approx\pi\left[3(a+b)-\sqrt{(3a+b)(a+3b)}\right]$.\\
        Para $a=2$, $b=1$: \ $L\approx9.6884$.
      \end{observacion}
    \end{column}
    \begin{column}{0.42\textwidth}
      \centering
      \begin{tikzpicture}
        \begin{axis}[bma, xmin=-2.4, xmax=2.4, ymin=-1.4, ymax=1.4, axis equal image, width=6.3cm, height=4.6cm, xtick={2}, xticklabels={$a$}, ytick={1}, yticklabels={$b$}]
          \addplot[curva, domain=0:360, samples=100] ({2*cos(x)},{sin(x)});
        \end{axis}
      \end{tikzpicture}

      \small Las integrales elípticas aparecen también en el periodo exacto del péndulo.
    \end{column}
  \end{columns}
\end{frame}

\begin{frame}{La función longitud de arco}
  \begin{definicion}
    Para una curva suave $\mathbf r(t)$, $t\ge\alpha$:\qquad
    $\displaystyle s(t)=\int_\alpha^t\sqrt{x'(u)^2+y'(u)^2}\du$.
  \end{definicion}
  Por el TFC1:\qquad $\dfrac{ds}{dt}=\sqrt{x'(t)^2+y'(t)^2}=|\mathbf v(t)|$\quad (la \emph{rapidez}).
  \begin{ejemplo}[7: espiral logarítmica]
    $x=e^t\cos t$, $y=e^t\sen t$: \ $x'^2+y'^2=2e^{2t}$, \ $s(t)=\displaystyle\int_0^t\sqrt2e^u\du=\sqrt2(e^t-1)$.\\[1mm]
    Despejando $t=\ln\!\left(1+\frac{s}{\sqrt2}\right)$ obtenemos una \resaltar{parametrización por longitud de arco}
    (rapidez constante $=1$), útil en física y geometría diferencial.
  \end{ejemplo}
  \begin{observacion}
    Cuando $t\to-\infty$ la espiral da infinitas vueltas alrededor del origen, pero $s(0)-s(-\infty)=\sqrt2$: ¡longitud finita!
  \end{observacion}
\end{frame}

\section{Curvas polares}

\begin{frame}{Longitud de arco en coordenadas polares}
  Una curva $r=f(\theta)$, $\alpha\le\theta\le\beta$, es paramétrica con parámetro $\theta$:
  \[ x=f(\theta)\cos\theta,\qquad y=f(\theta)\sen\theta. \]
  Entonces $x'=f'\cos\theta-f\sen\theta$, \ $y'=f'\sen\theta+f\cos\theta$, \ y
  \[ x'^2+y'^2=f(\theta)^2+f'(\theta)^2. \]
  \begin{columns}
    \begin{column}{0.55\textwidth}
      \begin{teorema}
        \[ L=\int_\alpha^\beta\sqrt{r^2+\left(\frac{dr}{d\theta}\right)^2}\dth. \]
      \end{teorema}
    \end{column}
    \begin{column}{0.45\textwidth}
      \centering
      \begin{tikzpicture}[scale=1.1]
        \draw[->] (0,0) -- (2.6,0);
        \draw (0,0) -- (35:2.2);
        \draw (0,0) -- (45:2.2);
        \draw[very thick, uniAzul] (35:2) arc (35:45:2);
        \draw[uniGranate, thick] (45:2) -- (45:2.4);
        \draw[uniVerde, thick] (35:2) -- (37:2.4) -- (45:2.4);
        \node[uniAzul] at (38:1.6) {\scriptsize $r\,d\theta$};
        \node[uniGranate] at (49:2.45) {\scriptsize $dr$};
        \node[uniVerde] at (35:2.65) {\scriptsize $ds$};
      \end{tikzpicture}

      \small $ds^2\approx(r\,d\theta)^2+(dr)^2$.
    \end{column}
  \end{columns}
\end{frame}

\begin{frame}{Ejemplos polares}
  \begin{columns}
    \begin{column}{0.62\textwidth}
      \begin{ejemplo}[8: cardioide $r=a(1+\cos\theta)$]
        \[ r^2+r'^2=a^2\big[(1+\cos\theta)^2+\sen^2\theta\big]=2a^2(1+\cos\theta)=4a^2\cos^2\tfrac\theta2. \]
        Por simetría:
        \[ L=2\int_0^{\pi}2a\cos\tfrac\theta2\dth=8a\Big[\sen\tfrac\theta2\Big]_0^\pi=8a. \]
      \end{ejemplo}
      \begin{ejemplo}[9: circunferencia $r=2a\cos\theta$]
        $r^2+r'^2=4a^2\cos^2\theta+4a^2\sen^2\theta=4a^2$, \ $0\le\theta\le\pi$:
        \quad $L=\int_0^\pi2a\dth=2\pi a$. \ (Se recorre una vez en $[0,\pi]$.)
      \end{ejemplo}
    \end{column}
    \begin{column}{0.38\textwidth}
      \centering
      \begin{tikzpicture}
        \begin{axis}[bma, xmin=-0.7, xmax=2.3, ymin=-1.5, ymax=1.5, axis equal image, width=5.4cm, height=5.4cm, xtick={2}, xticklabels={$2a$}, ytick=\empty]
          \addplot[curva2, domain=0:360, samples=200] ({(1+cos(x))*cos(x)},{(1+cos(x))*sin(x)});
          \addplot[curva, dashed, domain=0:180, samples=100] ({cos(x)*cos(x)},{cos(x)*sin(x)});
        \end{axis}
      \end{tikzpicture}

      \small Cardioide (granate) y circunferencia $r=\cos\theta$ (azul).
    \end{column}
  \end{columns}
\end{frame}

\begin{frame}{Ejemplo 10: espiral logarítmica}
  \begin{columns}
    \begin{column}{0.55\textwidth}
      \begin{ejemplo}
        $r=e^{\theta/5}$, \ $0\le\theta\le4\pi$. \ $r'=\frac15e^{\theta/5}$:
        \[ \sqrt{r^2+r'^2}=e^{\theta/5}\sqrt{1+\tfrac1{25}}=\frac{\sqrt{26}}{5}e^{\theta/5}. \]
        \[ L=\frac{\sqrt{26}}{5}\int_0^{4\pi}e^{\theta/5}\dth=\sqrt{26}\left(e^{4\pi/5}-1\right)\approx57.8. \]
      \end{ejemplo}
      \begin{observacion}
        La espiral logarítmica corta a todos los radios con el mismo ángulo. Aparece en conchas de nautilo,
        galaxias y en la trayectoria de algunos insectos hacia la luz.
      \end{observacion}
    \end{column}
    \begin{column}{0.45\textwidth}
      \centering
      \begin{tikzpicture}
        \begin{axis}[bma, xmin=-8, xmax=13.5, ymin=-9, ymax=6, axis equal image, width=6.5cm, height=6cm, xtick=\empty, ytick=\empty]
          \addplot[curva, domain=0:720, samples=300] ({exp(rad(x)/5)*cos(x)},{exp(rad(x)/5)*sin(x)});
        \end{axis}
      \end{tikzpicture}
    \end{column}
  \end{columns}
\end{frame}

\begin{frame}{Actividad colaborativa: Medir una ruta}
  \small
  \begin{activa}[Aprendizaje basado en la experiencia]
    \textbf{(a)} Calcule la longitud del sendero $y=\frac23x^{3/2}$, $0\le x\le3$. \textbf{(b)} Halle la longitud de la cardioide $r=1+\cos\theta$. \textbf{(c)} En una hoja de cálculo, aproxime (a) con poligonales de $n=4,8,16$ segmentos y observe la convergencia.
  \end{activa}
  \begin{itemize}
    \item \textbf{Organización:} equipos de 3 a 4 con roles rotativos (coordinador, redactor, verificador, portavoz); 12--15 minutos.
    \item \textbf{Producto:} Resultados exactos, tabla de poligonales y conclusión.
    \item \textbf{Retroalimentación:} el docente recorre los equipos, hace preguntas y no da la solución hasta la puesta en común.
  \end{itemize}
\end{frame}

\section{Ejercicios}

\begin{frame}{Ejercicios propuestos}
  \begin{ejercicio}
    \begin{enumerate}
      \item Longitud de $y=\dfrac{x^3}{6}+\dfrac{1}{2x}$, $1\le x\le2$. \ \small(Verifique que $1+y'^2$ es un cuadrado.)
      \item Longitud de $x=\dfrac{y^4}{8}+\dfrac{1}{4y^2}$, $1\le y\le2$.
      \item Longitud de $y=\ln(1-x^2)$, $0\le x\le\frac12$. \ \small(Fracciones parciales.)
      \item Longitud de un arco de la cicloide $x=a(t-\sen t)$, $y=a(1-\cos t)$.
      \item Longitud de $x=e^t\cos t$, $y=e^t\sen t$, $0\le t\le\pi$.
      \item Longitud de $x=t-\tanh t$, $y=\sech t$ (tractriz), $0\le t\le T$.
      \item Longitud total de la curva polar $r=a\sen^3(\theta/3)$.
      \item Longitud de la espiral de Arquímedes $r=\theta$, $0\le\theta\le2\pi$ (déjela planteada; se resolverá con sustitución trigonométrica).
      \item Longitud de $r=\dfrac{2}{1+\cos\theta}$, $0\le\theta\le\frac\pi2$ (parábola).
      \item Pruebe que la longitud de $y=\cosh x$ entre $0$ y $a$ es igual al área bajo la curva entre $0$ y $a$.
    \end{enumerate}
  \end{ejercicio}
\end{frame}

\begin{frame}{Metacognición y uso responsable de la IA}
  \small
  \begin{columns}[t]
    \begin{column}{0.5\textwidth}
      \begin{metacog}[cierre, 3 min]
        \begin{enumerate}
          \metacomun
          \item ¿Por qué a veces la integral de longitud no tiene primitiva elemental? ¿Cómo procedería entonces?
        \end{enumerate}
      \end{metacog}
    \end{column}
    \begin{column}{0.5\textwidth}
      \begin{iabox}
        Pida a una IA la longitud de la cardioide. Verifique que use $\theta\in[0,2\pi]$ y el valor absoluto adecuado.\\[2pt]
        \textit{Regla:} primero intente usted; luego use la herramienta para contrastar, nunca para sustituir su razonamiento. Anote el cambio en su bitácora.
      \end{iabox}
    \end{column}
  \end{columns}
\end{frame}

\begin{frame}{Resumen}
  \begin{itemize}
    \item Longitud $=$ supremo de longitudes de poligonales inscritas.
    \item Paramétrica: $L=\displaystyle\int_\alpha^\beta\sqrt{x'(t)^2+y'(t)^2}\dt$ (recorrida una vez).
    \item Gráfica: $L=\displaystyle\int_a^b\sqrt{1+f'(x)^2}\dx$.
    \item Polar: $L=\displaystyle\int_\alpha^\beta\sqrt{r^2+r'(\theta)^2}\dth$.
    \item $ds/dt=|\mathbf v|$; muchas longitudes llevan a integrales no elementales (elípticas).
  \end{itemize}
  \vspace{3mm}
  \textbf{Próxima sesión:} sustitución trigonométrica y primitivas de funciones hiperbólicas.
\end{frame}


\begin{frame}{Evaluación y trabajo posterior}
  \small
  \begin{evalbox}[Práctica del miércoles de la semana 5: Dirigida 3]
    \textbf{Cobertura:} semanas 3 y 4. \textbf{Naturaleza:} práctica dirigida en equipo, con retroalimentación inmediata.
  \end{evalbox}
  \begin{itemize}
    \item \criterios
    \item \declaracionia
    \item \textbf{Portafolio:} guarde el producto de la actividad, la bitácora de metacognición y los ejercicios propuestos de esta sesión; son evidencia de progreso y base de la retroalimentación.
  \end{itemize}
\end{frame}

\end{document}
